\chapter{Quantization}
\label{chap:quantization}

% Separate chapter approved by the author on 2026-11-10.
% Outline only: no experimental conclusions are asserted here.
% Sources for subsequent writing: results/v16, results/v17, results/v18;
% tranad_plus/campaigns/precision and tranad_plus/campaigns/integer.

\section{Quantization methods and experimental protocol}
\label{sec:quantization-protocol}
% Identify parent checkpoints, datasets, architectures, configurations and seeds.
% Distinguish quantization schemes as well as storage and execution formats.

\section{Quantized weight storage and floating-point calculations}
\label{sec:quantization-weight-storage}
% Explain weight quantization followed by dequantization for FP32 computation.

\section{Integer matrix operations with quantized weights and activations}
\label{sec:quantization-integer-operations}
% Describe actual integer matrix products and the remaining floating-point steps.
% Distinguish four-bit operand precision from native INT4 hardware instructions.

\section{Anomaly-detection performance}
\label{sec:quantization-detection}
% Compare with paired FP32 parents under explicit threshold and PA protocols.

\section{Storage and computational requirements}
\label{sec:quantization-resources}
% Separate parameter counts, stored bytes, operation estimates and measured runtime.

\section{Limitations and implications for onboard inference}
\label{sec:quantization-onboard}
% State which results depend on the quantization scheme and evaluated platform.
